\documentclass[12pt]{article}
\usepackage[utf8]{inputenc}
\usepackage[T1]{fontenc}
\usepackage[margin=1in]{geometry}
\usepackage{setspace}
\usepackage{xcolor}
\usepackage[colorlinks=true, urlcolor=blue!60!black, linkcolor=blue!60!black]{hyperref}
\setlength{\parindent}{0em}
\setlength{\parskip}{0.6em}
\renewcommand{\baselinestretch}{1.0}
\begin{document}
\begin{center}
{\LARGE \textit{Lochner v. New York}, 198 U.S. 45 (1905)\par}
\vspace{0.5em}
{\large Case Brief --- Two-Page Edition\par}
\end{center}

\section{Memory Jogger}
New York's ten-hour, sixty-hour cap on bakers' work violated the ``liberty of contract'' protected by due process, lacking a direct relation to health (\href{https://caselaw.findlaw.com/court/us-supreme-court/198/45.html}{FindLaw}).

\section{Detailed Case Facts}
Section~110 of New York's Bakeshop Act barred bakery employees from working more than sixty hours a week or ten hours a day, with no emergency exception; other sections regulated sanitation. Utica baker Joseph Lochner, a repeat offender, was indicted for permitting an employee to work more than sixty hours in a week.

\section{Procedural History}
Lochner was convicted and fined \$50. The Appellate Division (3--2) and the New York Court of Appeals (4--3) affirmed, sustaining the law as a health measure.

\section{Judicial Votes}
Reversed, 5--4. Peckham wrote for Fuller, Brewer, Brown, and McKenna. Harlan, joined by White and Day, dissented; Holmes dissented alone.

\section{Holding}
Due process protects liberty of contract, including the right to buy and sell labor. A law restricting labor contracts between competent adults must bear a direct relation to a legitimate police-power end; Section~110 was not a valid health law.

\section{Analysis of Opinions}
\textbf{Majority:} Bakers are competent adults, not wards of the state; bread's wholesomeness does not turn on bakers' hours; and baking is not an unusually unhealthy trade, unlike the mining in \textit{Holden v. Hardy}. Courts must ensure the police power is not a ``mere pretext,'' and the separate sanitary sections showed the hours cap was a labor law in disguise. \textbf{Harlan:} with ``room for debate and for an honest difference of opinion'' about bakers' health, the Court should defer. \textbf{Holmes:} the case was decided ``upon an economic theory which a large part of the country does not entertain,'' and liberty should not block a dominant opinion unless fundamental principles are plainly infringed.

\section{Examples: Future Applications}
\textbf{Same-Side:} a 48-hour weekly cap for adult retail clerks resting on general fatigue would fall. \textbf{Opposite-Side:} bakery ventilation and fire-exit rules directly protect health and would stand. \textbf{Fence-Sitter:} a ten-hour cap for intercity bus drivers, whose fatigue directly endangers the public.

\section{Critique}
Critics call \textit{Lochner} judicial enforcement of laissez-faire; Sunstein faults its treatment of common-law baselines as neutral, and the Court repudiated it in \textit{West Coast Hotel} (1937). Originalists reject substantive due process altogether, while revisionists like Bernstein defend the majority as principled.

\section{Key Quotations}
\begin{quote}
``General propositions do not decide concrete cases. The decision will depend on a judgment or intuition more subtle than any articulate major premise.''
\end{quote}
\noindent --- Justice Holmes, dissenting, 198 U.S. at 76.

\end{document}
